\documentclass[12pt]{article}
\usepackage[margin=0.6in,top=0.85in,bottom=0.55in]{geometry}
\usepackage{lmodern}
\usepackage{microtype}
\usepackage{fancyhdr}
\usepackage{titlesec}
\usepackage{enumitem}
\usepackage{lastpage}
\usepackage[hidelinks]{hyperref}

\setlength{\parindent}{0pt}
\setlength{\parskip}{0.22em}
\setlength{\headheight}{26pt}
\titleformat{\section}{\large\bfseries\scshape}{}{0pt}{}

\pagestyle{fancy}
\fancyhf{}
\fancyhead[C]{%
\footnotesize
Student Name: \hspace{2.0in} \quad Assignment ID: U1F1 \quad Page \thepage\ of \pageref{LastPage}
}
\renewcommand{\headrulewidth}{0.5pt}
\renewcommand{\footrulewidth}{0pt}

\newcommand{\answerbox}[2]{%
\vspace{0.12em}
\noindent\textbf{#1}\par
\vspace{0.04em}
\noindent\fbox{%
\begin{minipage}[t][#2]{\dimexpr\textwidth-2\fboxsep-2\fboxrule}
\rule{\textwidth}{0pt}
\vspace{#2}
\end{minipage}
}\par
\vspace{0.22em}
}

\newcommand{\evidencebox}[1]{%
\noindent\fbox{%
\begin{minipage}{\dimexpr\textwidth-2\fboxsep-2\fboxrule}
\vspace{0.3em}
#1
\vspace{0.3em}
\end{minipage}}
}

\begin{document}

\begin{center}
{\LARGE\bfseries Logic Dungeon I}\par\vspace{0.04em}
{\large\scshape The Hall of Inscriptions}\par\vspace{0.08em}
\rule{0.78\textwidth}{0.5pt}\par\vspace{0.12em}
\end{center}

\textbf{Purpose:} Apply the logic tools from Weeks 1--3 to one unfamiliar case.

\section*{Opening: The Twenty Adventurers}

Twenty adventurers enter the first underground passage toward \textit{Castrum Malae Ratiocinationis}. The Hall of Inscriptions contains three possible exits: the \textbf{Laurel Tunnel}, the \textbf{Echo Tunnel}, and the \textbf{Slate Tunnel}. Some inscriptions are true but prove less than travelers think. Others use unstable words or claims broader than the evidence.

Use only the supplied evidence. If a fact cannot be settled from it, say so. At a danger station, careful diagnosis or repair rescues the endangered adventurer.

\evidencebox{
\textbf{Tools available:} conclusion, stated reason, hidden assumption, follow-test; stable terms and shifting meanings; subject and predicate; quality and quantity; A/E/I/O forms; careful scope; the limits of logic.
}

\vspace{0.25em}
\evidencebox{
\textbf{Route marks:}
\begin{itemize}[leftmargin=1.4em,itemsep=0.1em,topsep=0.15em]
\item Laurel doors bear a leaf mark.
\item Echo doors bear a bell mark.
\item Slate doors bear a square mark.
\item The evidence at later stations will determine which marks are supported, disputed, or still unknown.
\end{itemize}
}

\vfill
\textit{Proceed carefully. Decoration, confidence, and repetition are not proof.}

\newpage

\section*{Part 1: The Gate of Applause --- Danger Station}

Above the Laurel entrance, an inscription reads:

\evidencebox{``Travelers have applauded the Laurel Tunnel for a hundred years. Therefore the Laurel Tunnel is the wisest route for our company.''}

\textbf{Part 1 Q1:} Identify the conclusion, stated reason, and hidden assumption. Decide whether the conclusion follows.

\answerbox{Part 1 Q1 ANSWER BOX}{2.45in}

\textbf{Part 1 Q2:} Rewrite the conclusion so it states only what the applause actually supports. Then name one kind of evidence needed to judge whether the tunnel is wise or safe.

\answerbox{Part 1 Q2 ANSWER BOX}{1.8in}

\newpage

\section*{Part 2: The Two-Voiced Guide --- Danger Station}

A guide stands beside the Echo entrance and argues:

\evidencebox{``The Echo passage is sound; the masons found no cracks in its first chamber. Every sound choice deserves our trust. Therefore choosing the Echo passage deserves our trust.''}

\textbf{Part 2 Q1:} Identify the word that changes meaning. Explain both meanings.

\answerbox{Part 2 Q1 ANSWER BOX}{1.65in}

\textbf{Part 2 Q2:} Rewrite the guide's argument so the key word has one stable meaning. Decide whether the rewritten conclusion follows from the supplied reason.

\answerbox{Part 2 Q2 ANSWER BOX}{2.35in}

\newpage

\section*{Part 3: The Beast Ledger --- Danger Station}

The expedition finds a ledger labeled \textbf{Sample from the first chamber only}. It records eight observed tunnel creatures:

\evidencebox{
\begin{tabular}{lll}
1. ash moth --- winged & 2. ash moth --- winged & 3. cave beetle --- winged \\
4. cave beetle --- not winged & 5. blind cricket --- not winged & 6. blind cricket --- not winged \\
7. glass ant --- winged & 8. glass ant --- not winged &
\end{tabular}
}

A clerk adds four claims:

\begin{enumerate}[label=\Alph*.,leftmargin=2em,itemsep=0.12em]
\item All observed ash moths are winged creatures.
\item No observed cave beetles are winged creatures.
\item Some observed glass ants are winged creatures.
\item Most creatures in the entire Hall are winged creatures.
\end{enumerate}

\textbf{Part 3 Q1:} For Claims A, B, and C, identify the standard form (A, E, I, or O), quality, and quantity. For Claim C, also identify its subject and predicate.

\answerbox{Part 3 Q1 ANSWER BOX}{2.15in}

\textbf{Part 3 Q2:} Which claims are contradicted by the ledger or reach beyond it? Explain, then repair one unsafe claim without inventing evidence.

\answerbox{Part 3 Q2 ANSWER BOX}{2.15in}

\newpage

\section*{Part 4: The Shortcut Inscription}

A Slate inscription reads:

\evidencebox{``Some doors marked with a square open toward the upper road.''}

The scout concludes: ``Therefore all square-marked doors open toward the upper road, and none lead downward.''

\textbf{Part 4 Q1:} State the inscription in standard form and identify its subject, predicate, quality, quantity, and A/E/I/O letter.

\answerbox{Part 4 Q1 ANSWER BOX}{2.0in}

\textbf{Part 4 Q2:} Explain exactly what follows and what does not. Repair the scout's conclusion so it stays within the inscription's scope.

\answerbox{Part 4 Q2 ANSWER BOX}{2.25in}

\newpage

\section*{Part 5: The Mapmaker's Oath}

The mapmaker says:

\evidencebox{``I copied every line carefully from an older map. Therefore this map gives the true condition of all three tunnels today.''}

No one in the company has inspected the tunnels beyond their first chambers. The older map's date and maker are unknown.

\textbf{Part 5 Q1:} Use the four-step test. Then distinguish what logic can decide here from what requires observation, testimony, or corroboration.

\answerbox{Part 5 Q1 ANSWER BOX}{2.7in}

\newpage

\section*{Part 6: The Door of Narrow Claims --- Rescue Station}

\textbf{Part 6 Q1:} Repair both claims below. Keep each claim's subject matter, but make its terms stable and its scope no broader than the supplied evidence.

\begin{enumerate}[label=\alph*.,leftmargin=2em,itemsep=0.15em]
\item ``The Echo passage is sound, so it is a sound decision.''
\item ``Most creatures in the Hall are winged.''
\end{enumerate}

\answerbox{Part 6 Q1 ANSWER BOX}{2.1in}

\newpage

\section*{Part 7: Final Decision --- The Supported Tunnel}

Before deciding, the company confirms two final observations:

\evidencebox{
\begin{itemize}[leftmargin=1.4em,itemsep=0.12em,topsep=0.15em]
\item One square-marked door was opened and observed to rise toward daylight.
\item No tunnel has been inspected all the way to its exit.
\end{itemize}
}

\textbf{Part 7 Q1:} Which route is best supported for the company's next cautious step: Laurel, Echo, or Slate? Cite at least two earlier stations. Explain one rejected route, state what remains unknown, and name the next observation that would most improve the decision.

\answerbox{Part 7 Q1 ANSWER BOX}{3.15in}

\textbf{Part 7 Q2:} Complete the reflection with specific details: ``Logic helped the expedition decide \rule{1.8in}{0.4pt}, but it could not by itself establish \rule{1.8in}{0.4pt}.''

\answerbox{Part 7 Q2 ANSWER BOX}{1.35in}

\end{document}
